\begin{frame}{Postopno odkrivanje v okolju align}
	% Naloga 3.2:
	% Za pravilno postopno odkrivanje morate na enem mestu uporabiti ukaz `only',
	% na dveh pa ukaz `onslide'.
	% Namig: prva vrstica enačbe se od druge podprosojnice naprej nadaljuje
	% z vsebino druge vrstice. To dosežete tako, da z ukazom `only' samo na prvi
	% podprosojnici postavite del prve vrstice od znaka & do znakov \\,
	% vključno z njimi:
	%     (a+b)^n \only<1>{&= \ldots \\}
	% Na poznejših podprosojnicah tako ni ne treh pik ne konca vrstice,
	% zato se za (a+b)^n izpiše naslednja vrstica.
	Dokaži \emph{binomsko formulo}: za vsaki realni števili \(a\) in \(b\) in za vsako naravno število \(n\) velja
	\begin{align*}
	(a+b)^n \only{&=\ldots \\}<1>
	\onslide<2->{&= (a+b) (a+b) \dots (a+b) \\}
	\onslide<3->{&= a^n + n a^{n-1} b + \dots + \binom{n}{k} a^{n-k} b^k + \dots + n a b^{n-1} + b^n \\}
	&= \sum_{k=0}^n \binom{n}{k} a^{n-k} b^k
	\end{align*}
\end{frame}
